% =====================================================================
%  Chapter 16 opener -- "Chapter at a Glance", three reading paths + the
%  ladder staircase, THE running example (the same 10-node tree T as
%  Ch.12/Ch.13/Ch.14/Ch.15, edges 1-2,1-3,2-4,2-5,3-6,3-7,5-8,7-9,9-10)
%  re-drawn with the depth/parent + tin/tout stamp overlay, the "three
%  machines, one question" cost chart, the "who owns what" comparison
%  (Ch.5 owns entry-only flattening vs Ch.9 owns RMQ-via-Cartesian vs
%  Ch.12 owns the DFS vs Ch.13 owns the fold vs THIS chapter owns
%  LCA+lifting vs Ch.19/20/22 own dynamic path updates), roadmap,
%  recognition guide, constraint->tool table, notation.
%  Exports fig:c11-tree and fig:c11-taxonomy.
%  All overlays below are hand-computed on T; every later section re-uses
%  exactly these numbers (depths d, parents par, the 2n-1 tour, first[],
%  the up table, tin/tout, distances) -- if a section disagrees, this
%  opener is wrong and the section is right; fix it here, never the
%  other way.
% =====================================================================

\begin{chaptertheorybox}{Chapter at a Glance}
Chapters~12--15 all rooted the tree and walked away with one number (or one
number per root).  This chapter asks the question they kept deferring: a
\emph{query} now names \emph{two arbitrary nodes} -- find their join point,
measure the route between them, fetch the $k$-th ancestor -- and there are
$Q$ up to $2\cdot10^{5}$ such queries on a static tree.  The single object
that answers all of them is the \term{lowest common ancestor}: the deepest
node that lies on both root-chains.

Everything here is \emph{one question, three machines}.  Machine~1 is the
honest walk: equalise the depths, then step both pointers up together --
$\bigO{h}$ per query, zero setup, and the baseline every proof argues
against (\secref{c11-idea}).  Machine~2 converts LCA into a range-minimum
query on the \(2n-1\) \term{tour of the DFS walk}: record every visit, keep
the first occurrence \(\code{first[]}\) of each node, and the LCA is the
shallowest node between the two firsts -- a \term{sparse table} answers
that in $\bigO{1}$ (\secref{c11-euler}).  Machine~3 precomputes
\(\code{up}[v][j]$, the \(2^j\)-th ancestor of every node, and climbs any
distance in $\bigO{\log n}$ bit-steps (\secref{c11-lifting}); the same
table, carrying a second payload \(\code{agg}[v][j]\), lifts \emph{any
associative aggregate} -- max, sum, xor, a count -- along a jump, which is
the merged old-Chapter-12 material of \secref{c11-aggregates}.  Machine~3
is what contests actually use: it upgrades \emph{every} earlier tool (the
\(k\)-th ancestor, distance-through-the-hinge, depth buckets for
CF~208E's cousins) for the price of $\bigO{n\log n}$ memory.

\textbf{You will learn to}
\begin{itemize}[nosep]
  \item state the ancestor chain property (\propref{c11-chain}) and the
        equalise-then-step walk, and read a tree's LCA table by hand;
        \hfill\lv{2}
  \item build the \(2n-1\) tour, \(\code{first[]}\), and prove that LCA is
        a range-minimum of depths over the tour span (\thmref{c11-eulerthm});
        \hfill\lv{3}
  \item answer that RMQ with a sparse table in \(\bigO{1}\) -- including
        the two-overlapping-blocks trick and the idempotence caveat that
        makes it legal for \(\min\) but never for sums; \hfill\lv{4}
  \item build \(\code{up}[v][j]\), climb \(k\) in \(\bigO{\log n}\), and
        write the descending-pairs LCA that \emph{returns}
        \(\code{par}[u]\) -- not \(u\); \hfill\lv{3}--\lv{4}
  \item carry an aggregate on each jump (\thmref{c11-agg}), split a path
        query at the hinge, descend to the first ancestor past a
        threshold, and walk to the \(k\)-th node \emph{on a route};
        \hfill\lv{4}--\lv{5}
  \item turn the hinge into measurements -- distances, on-path tests,
        per-depth buckets -- and combine everything on the Summit
        CF~1062E. \hfill\lv{5}
\end{itemize}

\textbf{Prerequisites.} Chapter~12 (\chref{c09}) for the rooted-view DFS
and the depth/parent measures -- the setup DFS here is \coderef{c09-dfs3}
verbatim.  Chapter~5 (\chref{c24}) for the entry-only \(\tin/\tout\)
flattening we reuse as an ancestry test (but \emph{not} for the tour: the
\(2n-1\) LCA tour is a different object, contrasted in 16.2).  Chapter~13
(\chref{c10}) for the subtree fold whose pair-harvest this chapter makes
queryable, and for the distance identity \(\dist(s,u)+\dist(u,v)=\dist(s,v)\)
iff \(u\) is on the route.  Chapter~19 (\chref{c13}) is \emph{not} needed:
the sparse table is built from scratch here.
\end{chaptertheorybox}

% ---------------- Reading paths + ladder position ----------------
\begin{chaptertheorybox}{How to Read This Chapter (three paths)}
Level badges \(\lv1\)--\(\lv6\): \(\lv2\) is computing an LCA by the naive
walk on a drawn tree; \(\lv3\) is either machine built correctly under time
pressure (a tour table, or an \(\code{up}\) table with the sentinel
discipline); \(\lv4\) needs \emph{two} ideas composed (sparse table over
the tour, or a hinge-split aggregate query); \(\lv5\) is a query that
reduces to LCA through a translation layer you must invent yourself (route
parity for CF~1304E, depth buckets for CF~208E, the removal analysis of
the Summit); \(\lv6\) is knowing which machine -- or which decomposition
from Ch.19--24 -- the constraint table is actually asking for.

\textbf{Walk-first path (\(\lv2\)--\(\lv3\), "I have never seen an LCA").}
Read \secref{c11-idea} end to end: ancestors, the chain property, the
hand-computed census \(\exref{c11-census}\), and the equalise-then-step
walk.  Then \secref{c11-lifting}: the \(\code{up}\) table and CSES~1687's
\(k\)-th ancestor are the smallest complete win in the chapter.  Only then
return to \secref{c11-euler} and see the tour as a \emph{reduction}.

\textbf{Lifting-first path (\(\lv3\)--\(\lv4\), "I need the workhorse").}
\secref{c11-idea} for notation, then \secref{c11-lifting} twice: once for
the table, once for the descending-pairs LCA and its \(u=v\) shortcut.
Then \secref{c11-aggregates} -- the second table \(\code{agg}[v][j]\), the
two accumulators, the hinge split -- and CF~1304E.  This path covers what
the majority of contest problems actually use.

\textbf{RMQ-first path (\(\lv4\), "I already know sparse tables").}
Go straight to \secref{c11-euler}: the \(2n-1\) tour, \(\code{first[]}\),
\thmref{c11-eulerthm}, and the two-overlapping-blocks query; CSES~1135 is
the case study.  Read the caveat in \secref{c11-aggregates} before reusing
the trick anywhere non-idempotent.  Backward ropes: Chapter~9's
\(\operatorname{RMQ}\to\) LCA-on-Cartesian-tree direction (\chref{c26}) is the mirror
image of what you just did.
\end{chaptertheorybox}

\begin{figure}[H]
\centering
\begin{tikzpicture}[diagram, x=9mm, y=7mm]
  \foreach \st/\lab in {1/lv1,2/lv2,3/lv3,4/lv4,5/lv5,6/lv6} {
    \fill[\lab!22, draw=\lab, line width=0.8pt] (\st-1,0) rectangle (\st,\st*0.52);
    \node[font=\scriptsize\bfseries, text=\lab] at (\st-0.5,\st*0.52-0.42) {L\st};}
  \node[annot, anchor=south west, font=\scriptsize]
    at (0.05,0.62) {16.1 chain property + naive walk (ex:c11-census)};
  \node[annot, anchor=south west, font=\scriptsize]
    at (1.05,1.14) {16.2 tour + first[] + thm:c11-eulerthm (CSES 1135)};
  \node[annot, anchor=south west, font=\scriptsize]
    at (2.05,1.66) {16.3 up table + climb + lifting-LCA (CSES 1687)};
  \node[annot, anchor=south west, font=\scriptsize, text=lv5]
    at (3.05,2.18) {16.4 agg on a jump, hinge split (CF 1304E)};
  \node[annot, anchor=south west, font=\scriptsize, text=lv5]
    at (4.05,2.70) {16.5 distances, on-path, buckets (CF 208E)};
  \node[annot, anchor=south west, font=\scriptsize, text=lv6]
    at (5.05,3.22) {Summit CF 1062E: RMQ on tin + LCA + tin/tout};
\end{tikzpicture}
\caption*{Chapter~16 on the six-step ladder: the census and the honest walk
sit at \(\lv2\); each machine used standalone is \(\lv3\); composing a
sparse table over the tour, or an aggregate over a jump, is \(\lv4\); the
four solved problems climb \(\lv3\)--\(\lv5\); the Summit CF~1062E is step~5
read from step~6.  Walk-first path \(=\) steps 2--3; full contest path runs
to~6.}
\end{figure}

% ---------------- The running example ----------------
\subsection*{The Running Example: One Ten-Node Tree, Three Machines}

\emph{One} tree does all the work -- the same \(T\) Chapters~12--15 used,
now wearing ancestor-query luggage.  Nodes \(1,\dots,10\), rooted at \(1\);
edges
\[
  1\!-\!2,\; 1\!-\!3,\; 2\!-\!4,\; 2\!-\!5,\; 3\!-\!6,\; 3\!-\!7,\; 5\!-\!8,\; 7\!-\!9,\; 9\!-\!10 .
\]
Child sets in that view: \(1\!\to\!\{2,3\}\), \(2\!\to\!\{4,5\}\),
\(3\!\to\!\{6,7\}\), \(5\!\to\!\{8\}\), \(7\!\to\!\{9\}\), \(9\!\to\!\{10\}\);
leaves \(\{4,6,8,10\}\); height \(4\), reached on the spine
\(1,2,5,8\) and the spine \(1,3,7,9,10\).  The tree is deliberately
\emph{thin}: depth \(4\) on \(n=10\) means the naive walk is already slow
enough to feel, and a chain-shaped worst case is only one re-root away.

\figref{c11-tree} stamps \(T\) with everything the machines consume.  Below
each node: \code{d} = depth (root at \(0\)), \code{p} = parent, with the
sentinel \(\code{par}[1]=0\) made visible; above-left/above-right:
\(\tin/\tout\) under the Ch.~5 entry-only convention with \(\tout\) = last
stamp inside the subtree.  These exact numbers are re-used, unchanged, in
every section: the \(2n-1\) tour of 16.2 is the DFS \emph{walk} over this
drawing, the \(\code{up}\) table of 16.3 is built from the \code{p} row,
and every distance in 16.5 is read off the \code{d} row.

\begin{figure}[H]
\centering
\resizebox{\linewidth}{!}{%
\begin{tikzpicture}[diagram, scale=0.90, transform shape,
    tnode/.style={circle, draw=primaryblue, line width=0.9pt, fill=primaryblue!8,
                  minimum size=8mm, font=\small\bfseries, text=primaryblue!60!black}]
  \node[font=\scriptsize\ttfamily, text=warningred!70!black] at (0,1.05) {$0$ (sentinel)};
  \draw[-latex, darkgray] (0,0.8) -- (0,0.32);
  \node[tnode] (n1) at (0,0)     {$1$};
  \node[tnode] (n2) at (-3,-1.6) {$2$};
  \node[tnode] (n3) at ( 3,-1.6) {$3$};
  \node[tnode] (n4) at (-5.2,-3.2) {$4$};
  \node[tnode] (n5) at (-1.6,-3.2) {$5$};
  \node[tnode] (n6) at ( 1.6,-3.2) {$6$};
  \node[tnode] (n7) at ( 5.2,-3.2) {$7$};
  \node[tnode] (n8) at (-1.6,-4.8) {$8$};
  \node[tnode] (n9) at ( 5.2,-4.8) {$9$};
  \node[tnode] (n10) at ( 5.2,-6.4) {$10$};
  \draw[tedge] (n1)--(n2) (n1)--(n3) (n2)--(n4) (n2)--(n5)
               (n3)--(n6) (n3)--(n7) (n5)--(n8) (n7)--(n9) (n9)--(n10);
  % depth / parent badges (hand-computed on T)
  \foreach \nd/\d/\p in {
      n1/0/0, n2/1/1, n3/1/1, n4/2/2, n5/2/2,
      n6/2/3, n7/2/3, n8/3/5, n9/3/7, n10/4/9} {
    \node[font=\tiny\ttfamily, text=warningred!75!black, anchor=north]
      at (\nd.south) {d=\d\ p=\p};}
  % tin / tout stamps (entry-only DFS, ascending children, tout = last stamp)
  \foreach \nd/\ti/\to in {
      n1/1/10, n2/2/5, n3/6/10, n4/3/3, n5/4/5,
      n6/7/7, n7/8/10, n8/5/5, n9/9/10, n10/10/10} {
    \node[font=\tiny\ttfamily, text=successgreen!50!black, anchor=south]
      at (\nd.north) {\ti/\to};}
  \node[annot, font=\scriptsize, anchor=north west, text=warningred!70!black]
    at (7.0,0.2) {\(\begin{array}{@{}l@{}}
      \text{depth } d=(0,1,1,2,2,2,2,3,3,4)\\
      \text{parent } p=(0,1,1,2,2,3,3,5,7,9)\\
      \text{tin: } 1\!:\!1,2\!:\!2,4\!:\!3,5\!:\!4,8\!:\!5,\\
      \quad 3\!:\!6,6\!:\!7,7\!:\!8,9\!:\!9,10\!:\!10\\
      \text{tout (last stamp): } 1\!:\!10,2\!:\!5,3\!:\!10,\\
      \quad 4\!:\!3,5\!:\!5,6\!:\!7,7\!:\!10,8\!:\!5,9\!:\!10,10\!:\!10
    \end{array}\)};
  \node[annot, font=\scriptsize, anchor=west, text=darkgray]
    at (-5.4,-6.9) {below node: depth, parent \quad above node: tin/tout \quad dashed spine 1-2-5-8 is one root-chain};
  \draw[dashed, accentorange, line width=0.9pt] (n1)--(n2)--(n5)--(n8);
\end{tikzpicture}%
}
\caption{The running tree \(T\) with the ancestor-query overlay. Below each node:
depth \(d\) (root at \(0\)) and parent \(p\); above each node: \(\tin/\tout\) from the
entry-only DFS of Chapter~5 (\chref{c24}), children visited in ascending label order,
\(\tout\) = the last stamp inside the subtree. The orange dashed spine
\(1\!-\!2\!-\!5\!-\!8\) is one root-chain, so \(\ancestors(8)=\{8,5,2,1\}\). Sanity
checks used across the chapter: node~10 is the unique depth-4 node; \(u\preceq v\) iff
\(\tin(u)\le\tin(v)\le\tout(u)\) --- e.g.\ \(7\preceq 10\) as \(8\le10\le10\), while
\(9\preceq 6\) fails as \(9>7\). Every tour, table and climb in 16.2--16.5 re-uses
exactly these stamps.}
\label{fig:c11-tree}
\end{figure}

Four angles, four hand-computed overlays -- each lives in the section named
on its row and never changes:

\begin{center}
\small
\renewcommand{\arraystretch}{1.25}
\begin{tabular}{@{}p{0.26\linewidth} p{0.46\linewidth} l@{}}
\toprule
\textbf{Angle on \(T\)} & \textbf{What comes out} & \textbf{Traced in} \\
\midrule
ancestor census & \(\lca(8,10)=1\), \(\lca(4,5)=2\), \(\lca(6,7)=3\), degenerate \(\lca(9,10)=9\) & \secref{c11-idea} \\
\(2n-1\) tour + RMQ & tour \(1,2,4,2,5,8,5,2,1,3,6,3,7,9,10,9,7,3,1\) (length 19), \(\code{first}[8]=5\), \(\code{first}[10]=14\), span min depth \(0\to\) LCA~1 & \secref{c11-euler} \\
\(\code{up}[v][j]\) table & \(\code{up}[10][2]=1\), \(\code{up}[8][1]=2\); climb for \(\lca(8,10)\) equalises \(10\to 9\), jumps \(2\!\to\!1\) region, returns \(\code{par}[2]=1\) & \secref{c11-lifting} \\
distances + buckets & \(\dist(8,10)=7\), \(\dist(4,8)=3\); depth-2 bucket \(\{3,4,7,8\}\) answers "cousins at distance~2" counts & \secref{c11-paths} \\
\bottomrule
\end{tabular}
\end{center}

% ---------------- Three machines, one question ----------------
\subsection*{Three Machines, One Question (the cost chart)}

All three answer the \emph{same} question -- \(\lca(u,v)\) on a fixed rooted
tree -- and the right one is chosen by the shape of the input, not by
taste.  Write \(h\) for the height and \(n\) for the node count.

\begin{center}
\small
\renewcommand{\arraystretch}{1.25}
\setlength{\tabcolsep}{4pt}
\begin{tabular}{@{}>{\raggedright\arraybackslash}p{0.18\linewidth} >{\raggedright\arraybackslash}p{0.13\linewidth} >{\raggedright\arraybackslash}p{0.14\linewidth} >{\raggedright\arraybackslash}p{0.13\linewidth} >{\raggedright\arraybackslash}p{0.32\linewidth}@{}}
\toprule
\textbf{Machine} & \textbf{Setup} & \textbf{Per query} & \textbf{Memory} & \textbf{When it wins} \\
\midrule
naive walk (16.1) & \(\bigO{1}\) (just \code{par}, \code{depth}) & \(\bigO{h}\) & \(\bigO{n}\) & one-off questions, shallow trees, proofs-by-baseline \\
Euler + RMQ (16.2) & \(\bigO{n\log n}\) & \(\bigO{1}\) & \(\bigO{n\log n}\) & pure LCA/distance queries only, nothing else needed from the table \\
binary lifting (16.3) & \(\bigO{n\log n}\) & \(\bigO{\log n}\) & \(\bigO{n\log n}\) & the default: LCA + \(k\)-th ancestor + aggregates + threshold descents \\
\bottomrule
\end{tabular}
\end{center}

\begin{insightbox}[One question is worth three machines]
The RMQ machine is \emph{faster} per query (\(\bigO{1}\)) but the lifting
machine is \emph{more useful}: the same \(\code{up}\) table also answers
the \(k\)-th ancestor, lifts max/sum/xor payloads, and descends to the
first ancestor past a threshold -- queries that have no RMQ translation at
all.  When the constraints fit both, contests pick lifting; the tour is
kept because the reduction LCA\(\to\)RMQ is one of the two or three great
problem-transformations in this book (Chapter~9 owns the reverse one).
\end{insightbox}

% ---------------- Taxonomy chart ----------------
\begin{figure}[H]
\centering
\begin{tikzpicture}[diagram, scale=0.95, transform shape,
    mnode/.style={circle, draw=#1, line width=0.9pt, fill=#1!8, minimum size=6.5mm,
                  font=\scriptsize\bfseries, text=#1!60!black}]
  \node[mnode=primaryblue] (a) at (0,0)     {walk};
  \node[annot, font=\scriptsize, anchor=west] at (1.15,0)  {equalise depths, step both up together; O(h); the baseline every proof argues against (16.1)};
  \node[mnode=successgreen] (b) at (0,-1)  {tour};
  \node[annot, font=\scriptsize, anchor=west] at (1.15,-1) {DFS walk of length 2n-1 + first[]; LCA = min-depth RMQ over the span (16.2)};
  \node[mnode=difforange]   (c) at (0,-2)  {up};
  \node[annot, font=\scriptsize, anchor=west] at (1.15,-2) {up[v][j] = 2\textasciicircum j-th ancestor; bit-climb in O(log n); the workhorse (16.3)};
  \node[mnode=warningred]   (d) at (0,-3)  {agg};
  \node[annot, font=\scriptsize, anchor=west] at (1.15,-3) {second table agg[v][j]: any associative payload on a jump; hinge split (16.4)};
  \node[mnode=softpurple]   (e) at (0,-4)  {hinge};
  \node[annot, font=\scriptsize, anchor=west] at (1.15,-4) {LCA as a measurement hub: distances, on-path tests, depth buckets, route positions (16.5, Summit)};
\end{tikzpicture}
\caption{The chapter's five-state taxonomy: rows 1--3 are the three
machines for the question "\(\lca(u,v)=?\)", row~4 is the same lifting
machine carrying a payload (the merged old-Chapter-12 topic), row~5 is
what the hinge \emph{measures}.  Every solved problem in the chapter sits
on exactly one row: CSES~1135 on row~2, CSES~1687 on row~3, CF~1304E on
row~4, CF~208E on row~5, the Summit CF~1062E across rows 2 and 5.}
\label{fig:c11-taxonomy}
\end{figure}

% ---------------- Who owns what ----------------
\subsection*{Who Owns What (so nothing is learned twice)}

\begin{center}
\small
\renewcommand{\arraystretch}{1.25}
\begin{tabular}{@{}p{0.30\linewidth} p{0.32\linewidth} p{0.32\linewidth}@{}}
\toprule
\textbf{Idea} & \textbf{Owner} & \textbf{This chapter's share} \\
\midrule
entry-only \(\tin/\tout\) flattening, subtree-as-interval & Ch.~5 (\chref{c24}) & reused as the ancestry test \(\propref{c11-interval}\); never re-derived \\
rooted DFS, depth/height/size measures & Ch.~12 (\chref{c09}) & \coderef{c09-dfs3} reused verbatim as the setup pass \\
subtree fold, pair-at-LCA harvest & Ch.~13 (\chref{c10}) & the harvest becomes \emph{queryable}; \exref{c11-census} is its worked basis \\
RMQ via Cartesian tree & Ch.~9 (\chref{c26}) & the reverse reduction; this chapter owns static sparse tables \\
LCA + lifting + lifted aggregates & \textbf{Ch.~16 (this)} & three machines, the hinge, and the merged old-Ch.~12 payload material \\
dynamic path updates / subtree updates & Ch.~20 (\chref{c15}), Ch.~22 (\chref{c16}) & \emph{not} here -- the moment values change between queries, hand off (HLD) \\
distance \emph{distributions} over all pairs & Ch.~24 (\chref{c18}) & counting over every pair is centroid territory; per-pair queries stay here \\
\bottomrule
\end{tabular}
\end{center}

% ---------------- Recognition guide ----------------
\begin{recognitionbox}[Pattern card -- ``the answer is a distance, route or k-th ancestor between named pairs'']
\reckey{a \emph{static} tree plus \(Q\) questions of the form: how far apart
are \(u\) and \(v\); is \(w\) on the path \(u\)--\(v\); who is the lowest
common boss of employees \(x,y\); the \(k\)-th ancestor / \(k\)-th node on
the route; the deepest ancestor whose value still passes a test; a count
inside the subtree at depth \(d\) from \(v\).}
\recuse{
  \item per-pair distance or LCA queries on a fixed tree -- any machine
        from this chapter; \secref{c11-euler} or \secref{c11-lifting};
  \item "\(k\)-th ancestor of \(v\)", or "first node above \(v\) with
        property \(P\)" -- \secref{c11-lifting} and the bit-descent of
        \secref{c11-aggregates};
  \item an extremum / sum / xor \emph{along a path}, values frozen --
        \(\code{agg}[v][j]\), \secref{c11-aggregates};
  \item "cousins at distance \(j\) from \(v\)" and similar depth-indexed
        subtree counts -- hinge + buckets, \secref{c11-paths} (CF~208E);
  \item walk \(a\to b\), detour through a fresh edge, parity/wiggle room
        -- route length via distances, \secref{c11-aggregates} (CF~1304E).
}
\recnotuse{
  \item edge or vertex values \emph{change} between queries -- that is
        segment trees + heavy-light, Chapters~20/22 (\chref{c15},
        \chref{c16}); nothing here survives an update;
  \item the answer must be reported for every \emph{root} -- rerooting,
        Chapter~15 (\chref{c20});
  \item counting pairs over the whole tree (a distance \emph{distribution})
        -- centroid decomposition, Chapter~24 (\chref{c18});
  \item a single BFS-solvable distance on an unweighted graph that happens
        to be a tree \emph{forest} with per-component roots -- check the
        input is really rooted first (the super-root trick, 16.3).
}
One-line test: \emph{do the queries name two nodes and ask what lies
\underline{between them} (or above one of them) on the unique route, while
the tree itself never changes?  Then this chapter's machine is the answer
structure.}
\end{recognitionbox}

% ---------------- Constraint -> tool ----------------
\subsection*{Reading the Constraints (which machine the limits want)}

\begin{center}
\small
\renewcommand{\arraystretch}{1.25}
\begin{tabular}{@{}p{0.30\linewidth} p{0.40\linewidth} p{0.24\linewidth}@{}}
\toprule
\textbf{Constraint signature} & \textbf{What it rules in/out} & \textbf{Tool} \\
\midrule
\(n,Q\le10^5\), tree is a chain in the worst case & \(\bigO{h}\) walk is \(\bigO{n}\) per query -- TLE; any \(\bigO{\log}\) machine fits & \secref{c11-lifting}, default \\
\(n\le10^5\), \(Q\le5\cdot10^4\), pure LCA/distance, nothing else & \(\bigO{1}\) RMQ query pays off; one sparse table, no second table & \secref{c11-euler}, CSES~1135 \\
queries also ask \(k\)-th ancestor / path extrema / threshold climbs & RMQ cannot express them; \(\code{up}\) + \(\code{agg}\) can & \secref{c11-lifting}, \secref{c11-aggregates} \\
"depth \(j\) subtree count" per query & buckets keyed by depth + \(\tin\)-range inside the subtree & \secref{c11-paths}, CF~208E \\
\(n,Q\le10^5\) but \(h\) small in test data & the walk \emph{passes} -- fragile; contests have anti-\(h\) tests & walk only as a checker \\
values updated between queries & everything here is static; the table goes stale & Ch.~20/22 hand-off \\
\bottomrule
\end{tabular}
\end{center}

% ---------------- Notation ----------------
\subsection*{Notation Used Only Here}

Root the tree at~\(1\) unless stated (every contest input in this chapter
either says so or is a forest given a super-root~\(0\); see 16.3's
sentinel discipline).  \(\code{par}[v]\) and \(\code{depth}(v)\) come from
one measuring DFS (\coderef{c09-dfs3}); \(\code{par}[\text{root}]=0\) and
vertex~\(0\) is a \term{sentinel}, never a real answer -- every table row
for~\(0\) is \(0\) and every loop must stop before returning it.
\(\ancestors(v)\) is the root-chain \(v, \code{par}[v], \dots, 1\), and
\(u\preceq v\) ("\(u\) is an ancestor of \(v\)", including \(u=v\)) is the
\term{ancestor order} whose chains \propref{c11-chain} exploits.  The
\(2n-1\) tour is \(\code{tour[]}\) with \(\code{first}[v]\) = index of
\(v\)'s first appearance and \(\code{dep}[\,]\) the depth alongside it
(16.2's convention: indices are \(\text{0-based}\)).  The sparse table is
\(\code{st}[j][i]\) = index of the minimum-depth entry among
\(\code{tour}[i..i+2^j-1]\).  \(\code{up}[v][j]\) = the \(2^j\)-th
ancestor of \(v\), \(\mathit{LOG}\) is the smallest integer with
\(2^{\mathit{LOG}}>n\) (\(\mathit{LOG}=4\) suffices for \(T\)); \(\code{agg}[v][j]\) is
the aggregate over the \(2^j\) vertices \emph{starting at \(v\) and going
up, excluding the landing vertex} -- the convention pinned in 16.4 and
never varied.  \(\dist(u,v)\) is the number of \emph{edges} on the route.

Every section of this chapter is a machine run on the \emph{same} tree on
the \emph{same} stamps: if a table here and a trace there disagree, this
opener is wrong and the section is right -- fix it here, never the other
way.

